Un sistema de detección de intrusiones o IDS (Intrusion Detection System) es un mecanismo de seguridad cuyo objetivo es monitorizar la actividad de un sistema o red para identificar distintos tipos de intrusiones, como por ejemplo accesos no autorizados, comportamientos sospechosos o acciones que puedan comprometer la confidencialidad, integridad y seguridad de la información. A diferencia de otros mecanismos preventivos, los IDS detectan posibles intrusiones y generan alertas para que sean analizadas por equipos de seguridad o administradores de la organización \cite{abdulganiyu2023systematic}.

Dependiendo del entorno en el que se encuentren, los IDS pueden clasificarse en sistemas basados en host (HIDS) y sistemas basados en red (NIDS). Los HIDS se instalan en dispositivos concretos y analizan la actividad local de estos, incluyendo procesos en ejecución, logs del sistema o conexiones tanto entrantes como salientes. Su principal ventaja es que permiten observar con detalle lo que ocurre en una máquina específica. Por otra parte, requieren ser instalados y proporcionarles un mantenimiento en cada equipo que se desea proteger. Por otro lado, los NIDS se sitúan en puntos estratégicos de la infraestructura donde se encuentre información sensible para capturar y analizar el tráfico que circula por la red. Un ejemplo de un lugar donde se puede situar es en la DMZ o zona desmilitarizada.

Los IDS también pueden ser clasificados según su forma de detección. En esta clasificación se encuentran los IDS basados en firmas y los IDS basados en anomalías. Los IDS basados en firmas comparan los datos obtenidos a partir de la actividad monitorizada con los que se encuentran en una base de datos de patrones. Cuando el tráfico coincide con una de las firmas contenidas en la base de datos, el sistema genera una alerta. Este método es eficaz frente a ataques conocidos, pero depende de que las firmas estén correctamente definidas y actualizadas. Esta es su mayor flaqueza, ya que tiene dificultades para detectar ataques nuevos o variantes nuevas de los ya existentes en la base de firmas. Por otro lado, los IDS basados en anomalías construyen un modelo del comportamiento normal del sistema o red y consideran sospechosas aquellas actividades que se desvían de dicho comportamiento. Esto permite detectar ataques nuevos, pero es mas posible que se generen falsos positivos si el tráfico benigno no es el habitual. Existen además los IDS híbridos, los cuales combinan ambas formas de detección. Estos aportan las ventajas de cada uno de los dos tipos, pero aumentan la complejidad del sistema \cite{abdulganiyu2023systematic,kunal2019mlids}.

En los últimos años, se ha estudiado la incorporación a estos de técnicas de inteligencia artificial y aprendizaje automático, especialmente en los IDS basados en anomalías. Estas técnicas permiten entrenar modelos capaces de aprender patrones a partir de datos y clasificar el tráfico en función de si es benigno o un ataque. Dependiendo del enfoque, cada muestra puede representar un paquete o un flujo de estos o una conexión que se describen mediante un conjunto de características, obteniendo en la salida del modelo el tipo de tráfico al que pertenece esta muestra \cite{kunal2019mlids}.

Sin embargo, la aplicación de técnicas de aprendizaje automático a los IDS también introduce algunas dificultades. Entre estas se encuentran la necesidad de disponer de conjuntos de datos que representen la realidad y ricos en contenido, el desbalanceo entre clases en estos conjuntos, la presencia de ruido en las muestras, el coste computacional del entrenamiento y la dificultad para interpretar la salida de ciertos modelos. Esto último ha provocado la aparición de sistemas de detección de intrusiones explicables, cuyo objetivo es no solamente clasificar el tráfico, sino proporcionar información sobre por la toma de decisiones a la hora de clasificar del modelo \cite{neupane2022xids}.
